\section{Conclusion}
\label{sec:conclusion}

We evaluated GPT-5.4-mini reconstructing 261 feature files from reverse-engineered stories. The main contribution is an auditable multi-domain evaluation that separates parser compliance from structural resemblance and shows, through exploratory defect auditing and perturbation, why embedding similarity cannot certify business correctness. For RQ1, the median was 0.853; the exact sign test did not show that it exceeded the descriptive 0.85 reference point ($p=0.209$, $\delta=0.054$). For RQ2, 255/258 files parsed (98.84\%, $p<0.001$). In the reviewed high-similarity sample, 25/60 files had a business-logic defect; labels had below-target agreement, consensus covered 15 files, and 45 were not relabelled. This project-stratified exploratory result does not estimate defect prevalence in the full corpus. GPT-5.4-mini is therefore supported as a drafting aid with human review of outcomes, conditions, values and criterion coverage, not as an autonomous acceptance-test generator. Future work should use independently written requirements, independent criterion-level review or execution against step definitions, and retrieval-augmented or open-source baselines.
